\section{Exam evening, 2026-11-12}
\label{sec:examevening}

\subsection*{Before}

Eat before 17:30. Bring pens, pencils, an eraser, a ruler for sketches, your student card, and the calculator the course allows. Be at 255 McCaul Street by 17:55. Do not study a new topic on Thursday. The useful thing to do that afternoon is to read the \emph{First moves} and \emph{Trap} boxes of every section once, which takes about twenty minutes.

\subsection*{The first five minutes}

Read all three questions before writing anything. Then, on the first page of rough work, one line per question: the law, the control volume or coordinates, the two assumptions that matter, and a rank from 1 to 3 by how sure you are that you can finish. Start with rank 1. This is the call from \secref{sec:call}{reading a question}, done on paper.

Look over the crib sheet once in those five minutes as well, so you know which equations it prints and which you will have to write yourself. If the paper prints the Navier--Stokes equations, the differential question is a reduction exercise; if it does not, write them from the box in \secref{sec:navierstokes}{Navier--Stokes} first.

\subsection*{The eighty-five minutes}

Roughly twenty-eight minutes a question. A question not moving after twenty minutes: write down the set-up you have (the control volume, the assumptions, the equation, its reduced form) and move to the next. The set-up carries most of a question's marks, and a finished set-up on all three beats a finished answer on one. Return with whatever time is left.

Leave the last five minutes for units. Every final number needs one, and so does every intermediate number that you box.

\subsection*{The write-up, every question}

\begin{enumerate}[nosep,leftmargin=*]
\item A sketch. The control volume as a dashed line, the inlets and outlets labelled 1 and 2, the coordinate axes with their positive directions, and the forces on the control surface as arrows. For a differential question: the geometry, the coordinates, and the walls with their velocities.
\item The assumptions, as a short list. Steady; incompressible; uniform velocity at each cut; gauge pressure; neglect gravity in $x$; fully developed; and so on.
\item The governing equation as printed or as written from memory, before any term is removed.
\item The reduced equation, with a reason beside each term removed: ``$\partial/\partial t = 0$, steady'' and ``$v = 0$, from continuity''.
\item The algebra, then the numbers, with units carried through.
\item One line of check: the sign of a force against intuition, a velocity at a wall against the no-slip condition, the dimensions of a Pi group.
\end{enumerate}

\subsection*{After}

That evening, write the three call lines from memory, as on the exam's first page, on paper or in a note that does not sync. Keep them out of \texttt{calls.md}, which is public on GitHub, until the course has posted solutions or returned the papers, because some students write the midterm later. Then type them in with \texttt{2026-11-12} at the start of each and the domain \texttt{mie312}, and settle them against the posted solutions or the marked paper. The \texttt{mie312} row under Domains in \texttt{readout.md} is the record of this test.
